\section*{Chapter \ref{ch:g12:populations-conservation} --- Populations and Conservation}

\begin{solution}{exo:g12:populations-conservation:1}
A \omterm{def:g12:populations-conservation:pop}{population} is the set of interbreeding individuals of one \omterm{def:g12:phylogeny-macroevolution:species}{species} in a
defined area (or potentially interbreeding, for many practical counts).
\omterm{def:g12:populations-conservation:K}{Carrying capacity} $K$ is the long-run size the environment can sustain given
resources and conditions; it is not a fixed magic number forever, but a useful
model ceiling.
\end{solution}

\begin{solution}{exo:g12:populations-conservation:2}
Exponential \omterm{def:g10:scientific-biology:properties}{growth} is $N(t)=N_0\mathrm{e}^{rt}$. It is a reasonable short-term
description when density dependence is weak --- resources abundant and crowding
effects negligible --- not when $N$ approaches $K$.
\end{solution}

\begin{solution}{exo:g12:populations-conservation:3}
In the logistic model, net \omterm{def:g12:populations-conservation:pop}{population} \omterm{def:g10:scientific-biology:properties}{growth} rate is zero at $N=K$ (births
balance deaths on average). Absolute increase $\mathrm{d}N/\mathrm{d}t$ is
maximal near $K/2$, the steep middle of the S-curve.
\end{solution}

\begin{solution}{exo:g12:populations-conservation:4}
\omterm{def:g12:populations-conservation:interactions}{Predation} is $+/-$ (predator gains, prey loses). \omterm{def:g12:populations-conservation:interactions}{Mutualism} is $+/+$ (both
benefit). \omterm{def:g12:populations-conservation:interactions}{Competition} is $-/-$ (both suffer from shared limiting resources).
Signs summarise \omterm{def:g11:evolution-mechanisms:fitness}{fitness} effects on the two parties.
\end{solution}

\begin{solution}{exo:g12:populations-conservation:5}
$1000=500\mathrm{e}^{r}$ implies $\mathrm{e}^{r}=2$, so
$r=\ln 2\approx 0.69\,\mathrm{year}^{-1}$. That is the continuous exponential
\omterm{def:g10:scientific-biology:properties}{growth} rate that doubles the \omterm{def:g12:populations-conservation:pop}{population} in one year under the model.
\end{solution}

\begin{solution}{exo:g12:populations-conservation:6}
The \omterm{prop:g12:populations-conservation:exclusion}{competitive exclusion} idea: complete competitors on a single limiting
resource cannot stably coexist; one excludes the other. Coexistence is possible
via \omterm{def:g10:ecosystems:habitat-niche}{niche} partitioning, disturbance, \omterm{def:g12:populations-conservation:interactions}{predation} that preferentially hits the
superior competitor, and other mechanisms that prevent one \omterm{def:g12:phylogeny-macroevolution:species}{species} from
monopolising all resources.
\end{solution}

\begin{solution}{exo:g12:populations-conservation:7}
Major anthropogenic threats include \omterm{def:g10:ecosystems:habitat-niche}{habitat} loss and fragmentation,
overexploitation, invasive \omterm{def:g12:phylogeny-macroevolution:species}{species}, pollution, and climate change. They often
interact (e.g.\ climate stress plus \omterm{def:g10:ecosystems:habitat-niche}{habitat} isolation).
\end{solution}

\begin{solution}{exo:g12:populations-conservation:8}
Small \omterm{def:g12:populations-conservation:pop}{populations} risk demographic stochasticity (chance birth/death runs),
inbreeding depression, strong \omterm{def:g11:evolution-mechanisms:drift}{genetic drift}, and local catastrophes with
little rescue by immigrants. Extinction probability rises even if mean \omterm{def:g10:scientific-biology:properties}{growth}
looks acceptable on paper.
\end{solution}

\begin{solution}{exo:g12:populations-conservation:9}
In situ conservation protects \omterm{def:g12:phylogeny-macroevolution:species}{species} in their natural \omterm{def:g10:ecosystems:habitat-niche}{habitat} (e.g.\ a
national park). Ex situ conservation maintains \omterm{def:g10:scientific-biology:organism}{organisms} outside the wild
(captive breeding, \omterm{def:g12:plant-biology:flower}{seed} banks, botanic gardens). Both can be \omterm{prop:g11:dna-replication:base-pairs}{complementary}
when wild \omterm{def:g10:ecosystems:habitat-niche}{habitat} is insufficient.
\end{solution}

\begin{solution}{exo:g12:populations-conservation:10}
Logistic \omterm{def:g10:scientific-biology:properties}{growth} is S-shaped: slow at first, steep mid-phase, flattening at
$K$. Raising $K$ (e.g.\ restoring wetland area or food resources) lifts the
ceiling so a larger sustainable \omterm{def:g12:populations-conservation:pop}{population} is possible if other threats are
controlled.
\end{solution}

\begin{solution}{exo:g12:populations-conservation:11}
In simple logistic models, surplus production available for harvest is
greatest near $K/2$; near $K$ net \omterm{def:g10:scientific-biology:properties}{growth} is small. Overharvest that drives $N$
too low can push the \omterm{def:g12:populations-conservation:pop}{population} into decline and collapse, especially with
Allee effects or environmental shocks.
\end{solution}

\begin{solution}{exo:g12:populations-conservation:12}
Protect core forest \omterm{def:g10:ecosystems:habitat-niche}{habitat}; build wildlife crossings or corridors across
roads so subpopulations connect; engage land users with incentives and
enforcement; monitor \omterm{def:g12:populations-conservation:pop}{populations} and \omterm{def:g10:ecosystems:habitat-niche}{habitat} so actions can be adjusted. A
single reserve without connectivity is often insufficient.
\end{solution}

\begin{solution}{pb:g12:populations-conservation:1}
\textbf{Part I.} (1)~$\mathrm{d}N/\mathrm{d}t=rN(1-N/2000)$.
(2)~Logistic rise from $200$ toward $2000$. (3)~No --- effective ceiling may
be set by \omterm{def:g12:populations-conservation:interactions}{predation}, not \omterm{def:g10:ecosystems:habitat-niche}{habitat} $K$ alone. (4)~$N=200\mathrm{e}^{1}
\approx 200\times 2.7=540$ under pure exponential for one time unit.

\textbf{Part II.} (5)~\omterm{def:g12:populations-conservation:interactions}{Predation}; intraspecific \omterm{def:g12:populations-conservation:interactions}{competition}.
(6)~E.g.\ seed-dispersal \omterm{def:g12:populations-conservation:interactions}{mutualism}; partner loss reduces food or \omterm{def:g10:ecosystems:habitat-niche}{habitat}.
(7)~Trophic cascades, release of other prey, vegetation change.

\textbf{Part III.} (8)~Invasives, \omterm{def:g10:ecosystems:habitat-niche}{habitat} loss (and possibly climate).
(9)~Often remove rats first (immediate mortality), protect \omterm{def:g10:ecosystems:habitat-niche}{habitat}, use
captive breeding if wild $N$ is critical. (10)~Undeveloped coastal strip
linking beaches. (11)~Counts, nesting success, rat indices, \omterm{def:g10:ecosystems:habitat-niche}{habitat} area.
(12)~Integrated narrative linking demography, threats and actions as asked.
\end{solution}
